% !TeX root = ../final-report.tex
% TODO(team): Recheck fixture and metric definitions if the study is rerun.
\section{Experimental Methodology}
\label{sec:method}

\paragraph{Experimental setup.}
The testbed consists of three Controllers, three storage Peers, and
independent viewers running on one host. Each run starts fresh Controller
and Peer processes. Viewers download segments over HTTP, verify their
hashes, and consume them in order at the rate specified by their media
durations. This consumption is simulated, while segment transfers and
bandwidth contention occur in the running system.
Table~\ref{tab:fixture} lists the cluster and client settings.

\paragraph{Workload.}
Each viewer requests the same 18\,s video, divided into 9 segments of 128\,KiB each.
Low demand consists of 2 viewers, while high demand consists of 6 viewers. 
Under high demand, arrivals follow 2 patterns.
With simultaneous arrivals, all 6 viewers start near 0\,s. 
With staggered arrivals, viewers start in pairs near 0, 4, and 12\,s. 

\begin{table}[ht]
\centering\small
\caption{Playback and storage configuration.}
\label{tab:fixture}
\begin{tabular}{@{}ll@{}}
\toprule
Storage cluster & 3 Controllers, 3 Peers \\
Source upload limits & 96, 144, 192 KiB/s \\
Content & 9 segments of 128 KiB and 2 s each \\
Initial placement & 2 copies per segment, fixed holders \\
Client prefetch & 3-segment window, 2 downloads at once \\
Startup buffer & 4 s of contiguous media \\
Timeouts & 8 s per attempt, 42 s per session \\
\bottomrule
\end{tabular}
\end{table}

\paragraph{Comparison design.}
Within each workload and arrival pattern, 
policies use the same content, initial replica placement, round-robin source selection, and fault schedule.
Policies are tested in a randomized order to reduce bias from changes in host load and cache state over time.
The main playback and repair comparisons cover 3 failure contexts, 
where the failed Peer's upload limit is set to 96, 144, or 192 KiB/s respectively.
The second-holder failure experiments use 2 of these contexts.

\paragraph{Fault schedule and recovery verification.}
The first Peer failure occurs 5\,s after the run starts.
The main observation window ends at 50\,s.
Replica recovery is checked approximately every 2\,s.
Essentially, full recovery requires every affected segment to have 2 reachable,
hash-valid copies with committed registrations.
Based on this, we count a run as meeting the recovery target if full recovery is verified within 40\,s of the first Peer failure.
The main study comprises 81 runs, including 72 performance runs and 9 Controller failure trials.

\paragraph{Metrics.}
To capture both playback performance and replica recovery, 
we measure session completion, startup delay, playback stall, replica recovery time, repair traffic, and one-copy exposure.

A session completes when it consumes all 18\,s of media.
For each session, startup delay measures the time from session start until media consumption begins. 
Once consumption begins, any time spent waiting for unavailable segments counts as playback stall.
To compare the effects of failure and repair, 
we accumulate observed stall from the first Peer failure until each session ends, 
then average over all scheduled viewers. 
For the no-failure control, accumulation begins at the corresponding 5-second mark.

Replica recovery time is measured from the first Peer failure
to the first check confirming full recovery.
Repair traffic counts source HTTP response payload bytes,
including repeated transfers.
When a segment has only one verified, reachable, registered copy,
failure of its remaining holder can make it unavailable.
Therefore, we use one-copy exposure to quantify the amount of data in this vulnerable
state and how long it remains there.
Let $b_s$ be the size of affected segment $s$ in MiB,
and let $n_s(t_i)$ be its verified, reachable, registered copy count
at check $i$.
The sampled exposure is

\begin{equation}
  E_{\mathrm{one}}
  = \sum_i (t_{i+1}-t_i)
    \sum_s b_s\,\mathbf{1}[n_s(t_i)=1],
\end{equation}

where the intervals span the period from the first failure
to the observation horizon.
Exposure is reported in MiB$\cdot$s.

